La presente memoria se organiza en seis capítulos cuyo contenido se describe a continuación:

El Capítulo 1 introduce el contexto y la motivación que justifican la realización del proyecto, expone los principales objetivos técnicos y pedagógicos planteados, y proporciona una visión general de la estructura del documento.

El Capítulo 2 recoge el estado del arte, analizando las tecnologías, metodologías y plataformas existentes más relevantes para el proyecto, incluyendo el uso de la IA en educación, los motores de sandboxing y los entornos interactivos de programación.

El Capítulo 3 presenta el análisis y diseño del sistema, describiendo la metodología de trabajo seguida, los requisitos funcionales y no funcionales, los casos de uso, el modelo de datos y la arquitectura de componentes.

El Capítulo 4 detalla el proceso de desarrollo e implementación, describiendo el entorno de trabajo, las tecnologías empleadas y las decisiones de diseño más relevantes adoptadas durante la construcción del sistema.

El Capítulo 5 recoge las pruebas realizadas sobre el sistema, tanto las pruebas de seguridad del sandbox como la validación pedagógica con usuarios finales.

El Capítulo 6 expone las conclusiones del proyecto y propone líneas de trabajo
futuro.